The problem of studying rare events is central to many areas of computer simulations. We recently proposed an approach to solving this problem that passes through the computation of the committor function, showing how it can be iteratively computed in a variational way while efficiently sampling the transition state ensemble. 
Here, we greatly ameliorate this procedure by combining it with a metadynamics-like enhanced sampling approach in which a logarithmic function of the committor is used as a collective variable. 
This procedure leads to an accurate sampling of the free energy surface in which transition states and metastable basins are studied with the same thoroughness. 
We show that our approach can be used in cases with the possibility of competing reactive paths and metastable intermediates. 
In addition, we demonstrate how physical insights can be obtained from the optimized committor model and the sampled data, thus providing a full characterization of the rare event under study. 